\begin{proposition}[The global meridians have only the prescribed toric direction]\hypertarget{tag-0311}{}\label{push:src:B.4.31}
\StatementRecord{push:src:B.4.31}{proved}{0311}{\SeriesRef{push:src:B.4.4}, \SeriesRef{push:src:B.4.19}, \SeriesRef{push:src:B.4.9}; \citet[Theorems~1.6 and 1.10, pp.~215--216, including the noncompact extension after Theorem~1.6]{BM97}}

Suppose $g:\mathcal Z\to\Delta_r$ is any proper complex-analytic
filling of a finite parameter cover of the six-sphere
family. Retain its reduced main component, normalize and resolve
it, extend its intrinsically marked map to $\PP^1$ by \ref{push:src:B.4.4},
resolve the graph, and resolve the support of its central divisor. Denote
the resulting smooth proper modification by $\widetilde{\mathcal Z}$,
and its holomorphic marked map by $h$. All modifications can be
chosen to be isomorphisms over $r\ne0$. The central divisor has normal-crossings support and positive integral
component multiplicities.

There is a finite subset $\Sigma\subset\PP^1$, containing
$0,1,\infty$, with the following property. For every
$q\in B^\circ:=\PP^1\setminus\Sigma$, every local angular
vanishing circle of the $r$-central normal-crossings divisor,
viewed in the homology of the nearby ordinary two-torus, lies in
\begin{equation}
             \Z\widehat\delta\subset H_1(X_{q,r};\Z).
\label{pub:I:eq:B-4-31-1}
\end{equation}
The same conclusion holds for the image of the entire local
angular torus at a higher crossing over $B^\circ$.
Here $\widehat\delta$ is the primitive cycle in the fixed
period marking. The conclusion is unaffected by any further
finite cover of the degeneration parameter.


\end{proposition}


\begin{proof}
\begin{enumerate}[label=\arabic*.,leftmargin=*]
\item \emph{The global model and the finite exceptional set.}
The punctured total family is connected and smooth, hence
irreducible. Its closure is the unique reduced main component
dominating the parameter disc. Retain this component, normalize
it, and resolve it properly, without changing its smooth
punctured part. The resulting total space is a smooth connected
fourfold, proper over the disc. The parameter is a nonzero
divisor, so its map to the disc is flat. The torsion-free coherent bi-anticanonical pencil and labelled
critical-value argument of \ref{push:src:B.4.4} apply to this smooth model.
Its marked graph is therefore a proper modification.
Further analytic resolution and principalization of the central
ideal give the stated smooth
model and normal-crossings support; the complex-analytic scope,
including locally finite resolution, is explicit in
\citet[Theorems~1.6 and 1.10, pp.~215--216, including the noncompact extension after Theorem~1.6]{BM97}. Properness over a smaller parameter disc
allows us to work over a compact central fibre throughout.
The modifications have their centres over $r=0$, so neither
the punctured family nor its marked map changes.
Flatness over $r$ persists because $r$ is a nonzero divisor in
the smooth local rings; all central multiplicities are kept in the
local monomial equations.

\item For possible self-crossings use the finite closed spaces of
ordered local branch intersections, with each local component
normalized. Locally these are the coordinate intersections of
the central branches. In particular, the double-branch spaces
are compact smooth surfaces and the triple-branch spaces are
compact smooth curves, even when two branches belong to one
global component. There are finitely many such components.

\item Add to $\Sigma$ the images of all branch-intersection components
on which $h$ is constant and the critical values of $h$ on every
remaining branch-intersection component, including the
normalizations of the single branches. These sets are finite.
Indeed the maps are proper; on a component mapping nontrivially
to the curve $\PP^1$ their critical image is a proper analytic
subset of that curve. Zero-dimensional deeper strata contribute
only finitely many further values.

\item Thus, over $B^\circ$, the closed double-branch spaces form
proper smooth families of compact curves. The connected
components of their fibres form finite covering spaces of
$B^\circ$: this follows either from proper submersion
trivializations, or from their finite Stein factorizations.
Removing the finitely many triple points from one connected
smooth curve leaves it connected. Branch orderings, central
multiplicities, and these connected fibre components therefore
constitute finite discrete data over $B^\circ$.

\item \emph{The local angular circles, including multiplicities.}
At a point of an open double-branch stratum, the submersion
condition for $h$ gives coordinates in which $q=h$ is a
tangential coordinate and
\begin{equation}
                 r=x^a y^b,\qquad a,b>0.
\label{pub:I:eq:B-4-31-2}
\end{equation}
A holomorphic unit has been absorbed in one coordinate on this
small chart. Put $d=\gcd(a,b)$. Each connected component of the
local nonzero fibre has the primitive angular circle
\begin{equation}
 \theta\longmapsto
 \bigl(x_0e^{i(b/d)\theta},\,
       y_0e^{-i(a/d)\theta}\bigr),
       \qquad 0\le\theta\le2\pi,
\label{pub:I:eq:B-4-31-3}
\end{equation}
where $x_0^a y_0^b=r$. There are at most $d$ choices of the
connected monomial-fibre component. Interchanging the ordered
branches can reverse the circle's orientation. These are only
finite choices.

\item Move the chosen base point along a path in the same open
double-curve fibre. Local product charts and the normal line
bundles transport the annular circle continuously. Changes of
normal frame rotate its parameter and do not add a tangential
one-cycle. Equivalently, this is the fibre circle of the normal
circle bundle along that path. If a based loop is used, changing
the path changes it by conjugation, which does not affect its
class in $H_1$. Consequently, for fixed small nonzero $r$ and
fixed $q$, each of the finitely many branch-component data
determines only finitely many homology classes, including signs.

\item The same local formula proves monodromy invariance in the
$r$ direction. Continuing around $r=0$ permutes the finitely
many connected components of $x^a y^b=r$; after a positive
power it carries \eqref{pub:I:eq:B-4-31-3} to the same oriented circle class.
In the marking, after a finite parameter cover the
$r$-monodromy is
\begin{equation}
 T_r=I-mE,\qquad
 E=\widehat\delta\otimes\gamma,\qquad m\in\Z_{>0},
 \quad E^2=0.
\label{pub:I:eq:B-4-31-4}
\end{equation}
For example, $w=r^6$ gives $m=12$ with the positive-loop
convention used in \ref{push:src:B.4.9}. If $T_r^k v=v$ for one of the
angular classes $v$, then $kmEv=0$. The homology lattice is
torsion-free, so
\begin{equation}
 v\in\ker E
    =\Z\widehat u+\Z\widehat w+\Z\widehat\delta.
\label{pub:I:eq:B-4-31-5}
\end{equation}

\item \emph{Why global continuation gives a finite orbit.}
Choose a base point $q_*\in B^\circ$. Let $\mathscr S$ be the
finite set of angular classes furnished at $q_*$ by all the
connected double-curve fibres, their ordered branches, their
monomial-fibre components, and both orientations.
Take generators of the finitely generated group
$\pi_1(B^\circ,q_*)$. A generator can be lifted to a path in
the appropriate family of double curves; it ends in one of
the finitely many connected fibre components. Avoid the
triple points during this lift. Properness over the compact
generator paths and finitely many product charts give one
sufficiently small $r_*$ for which all their annular circles
can be transported in the smooth fibres
$X_{q,r_*}$.

\item The transported circles are a continuous family of cycles,
so they realize the Gauss--Manin transport of $H_1$.
At the end, moving the base point within the resulting
connected punctured double curve does not change its $H_1$
class, by the previous paragraph. Hence each generator
permutes $\mathscr S$. The same holds for all words in the
generators. This proves a finite orbit in the marked lattice.
Its projection to
\begin{equation}
 \ker E/\Z\widehat\delta
       \simeq \Z\widehat u+\Z\widehat w
\label{pub:I:eq:B-4-31-6}
\end{equation}
therefore has finite orbit as well.

\item \emph{The exact elliptic-monodromy calculation.}
In the fixed marking the two finite peripheral
matrices are
\begin{equation}
 A_1=\begin{pmatrix}
 1&0&0&0\\6&0&1&0\\-6&-1&-1&0\\-2&1&0&1
 \end{pmatrix},\qquad
 A_2=\begin{pmatrix}
 1&0&0&0\\0&0&-1&0\\-6&1&0&0\\3&0&1&1
 \end{pmatrix}.
\label{pub:I:eq:B-4-31-7}
\end{equation}
They fix $\widehat\delta$ and $\gamma$, commute with $E$,
and induce on \eqref{pub:I:eq:B-4-31-6}
\[
 M_1=\begin{pmatrix}0&1\\-1&-1\end{pmatrix},\qquad
 M_2=\begin{pmatrix}0&-1\\1&0\end{pmatrix}.
\]
Deleting additional ordinary values from the elliptic base
does not decrease this monodromy image: the map on fundamental
groups to $\PP^1\setminus\{0,1,\infty\}$ is surjective.
The image therefore contains
\begin{equation}
 U=M_1M_2=\begin{pmatrix}1&0\\-1&1\end{pmatrix},\qquad
 V=M_2UM_2^{-1}=\begin{pmatrix}1&1\\0&1\end{pmatrix}.
\label{pub:I:eq:B-4-31-8}
\end{equation}
If a vector $\bar v$ has finite orbit, some positive powers
$U^k$ and $V^\ell$ both fix it. Their fixed lines are,
respectively, $\Q(0,1)$ and $\Q(1,0)$, whose intersection
is zero. Thus $\bar v=0$. Applying this to the projected
classes in $\mathscr S$ proves that each angular class lies
in $\Q\widehat\delta$. Since $\widehat\delta$ is a member
of an integral basis of $H_1(T^4;\Z)$, it is primitive;
an integral class in this rational line belongs to
$\Z\widehat\delta$. This proves \eqref{pub:I:eq:B-4-31-1}.

\item At a crossing with local equation
$r=\prod_{j=1}^k x_j^{a_j}$, the angular torus has
rational first homology
\[
 \ker\bigl((a_1,\ldots,a_k):\Q^k\longrightarrow\Q\bigr).
\]
It is spanned by the pairwise angular vectors
$a_j e_i-a_i e_j$. Each is a rational multiple of the
primitive circle of the corresponding double branch,
which can be moved to an open double-curve point by
keeping the other coordinates small and nonzero.
The same calculation therefore places the entire angular
image in $\Q\widehat\delta$, and primitivity again gives
the integral conclusion. This also handles self-crossings,
because the ordered local branches were retained throughout.
\end{enumerate}
\end{proof}